\documentclass[11pt]{article}
\usepackage[T1]{fontenc}
\usepackage{amsmath,amssymb,amsthm}
\usepackage[margin=1.12in]{geometry}
\usepackage[colorlinks=true,allcolors=blue,
  pdftitle={Nonzero sphere moments from nonordinary cyclic covers},
  pdfauthor={Roy van Rijn}]{hyperref}
\newtheorem{theorem}{Theorem}
\newtheorem{lemma}[theorem]{Lemma}
\newtheorem{proposition}[theorem]{Proposition}
\newcommand{\A}{\mathcal A}
\newcommand{\CT}{\operatorname{CT}}
\newcommand{\Q}{\mathbb Q}
\newcommand{\F}{\mathbb F}
\title{Nonzero sphere moments from nonordinary cyclic covers}
\author{Roy van Rijn}
\date{}

\begin{document}
\maketitle

\begin{abstract}
A sphere moment can detect Frobenius on one coherent line of a cyclic
cover even when its primitive Jacobian factor is nonordinary. For a fixed
prime \(p>3\) congruent to \(2\) modulo \(3\), we construct connected covers
of unbounded order \(N=(p^r-1)/(p-1)\), with odd \(r\ge3\), whose
primitive factors are nonordinary at \(p\) and are not potentially CM.
The sphere phases have least degree \(N\). For compatible lower
terms integral at \(p\), their moments of orders \((p^{kr}-1)/N\),
\(k\ge1\), multiplied by \(p^{kr}\), all reduce to one. At \(p=5\)
we give an exact primitive \(p\)-rank formula and a limiting rank
fraction of one half. In a second family, of order \(n=7M\) with
\(M\equiv5\pmod{14}\), the primitive factor is nonordinary at every
prime \(p\not\equiv1\pmod n\) with \(p\nmid n\) and \(c,1-c\)
units at \(p\). An orbit
criterion classifies the prime classes where the normalized moment
residue is a unit for orders \((p^{kd}-1)/n\), where \(d\) is the
order of \(p\) modulo \(n\) and \(k\ge1\). At least five classes
certify nonzero moments at nonordinary primes. For \(n=35\) the
nonordinary certifying classes have natural density \(7/24\).
\end{abstract}

\section{Introduction}

Consider the complex affine sphere \(2xy+t^2=1\). Substituting
\(y=(1-t^2)/(2x)\) sends a regular function to a Laurent polynomial in
\(x,t\). The normalized sphere moment is obtained by taking the constant
term in \(x\) and averaging over \(-1\le t\le1\). This defines a rational
functional \(\A\) over every characteristic-zero field.

Take a polynomial \(g\) of total degree less than \(31\) with rational
coefficients integral at \(5\). On the sphere set
\[
     f=t^{29}(x^2+12y^2-4t^2)+g.
\]
Its least sphere degree is \(31\), and both \(125\A(f^4)\) and
\(5^6\A(f^{504})\) reduce to one modulo five.
These and infinitely many further congruences come from one Frobenius line
on a cyclic cover of order \(31\). Its primitive Jacobian factor has
dimension \(60\) and \(5\)-rank \(42\), and a quotient of dimension
\(30\) has totally toric reduction at \(3\).

If the highest possible constant-term height in a power is \(p^r-1\),
its denominator \(p^r\) is isolated by reduction. In a cyclic cover
associated with the leading form, the coefficient that remains can be
the return of the \(r\)-fold Frobenius on a single line in coherent
cohomology. Ordinarity suffices to make such a return invertible.
Our examples show how to obtain an invertible return on the selected
line together with nonordinary reduction. The finer character
decomposition places the selected line and the failure of ordinarity
in separate summands.

In the first construction, the certifying prime and the first certified
moment order stay fixed as the conductor grows. At \(p=5\) we count the surviving
Frobenius orbits and show that the primitive \(p\)-rank is asymptotic to
half its dimension. The second construction classifies the primes
where the chosen return certifies infinitely many moment orders for a
fixed cover and phase. In both constructions,
a tame degeneration at a second prime rules out potential complex
multiplication on the primitive factor.

The fractional-part dimensions and Frobenius coefficient formulas for
tame cyclic covers belong to the Hasse--Witt theory of Bouw
\cite[\S\S2, 4--5]{bouw}. His generic rank theorem
\cite[Theorem~6.1]{bouw} attains the character-orbit bound for a
generic \(d\)-cyclic cover with \(b\) branch points when
\(p\ge d(b-3)\). Here we use the character calculation on explicit
covers and prove the selected coefficient is a unit in the specified
fixed-prime and fixed-cover regimes. For the three-branch quotient in
the second construction, the orbit condition agrees with the
Shimura--Taniyama slope calculation \cite[Theorem~3.2]{lmpt}. Here
we extract its exact residue in a sphere moment and use all seven
quotients to detect nonordinarity of the ambient primitive factor.
Ordinary reduction means that the \(p\)-rank equals the dimension.

\section{Sphere moments and cyclic characters}

For \(s=1,2\) assign weight \(j+si\) to the literal Laurent monomial
\(x^it^j\), and put \(z=x/t^s\). The weight-\(n\) part of a phase has
the form \(t^n B(z)\), where \(B\) is a Laurent polynomial. The
sphere functional on Laurent monomials is
\[
  \A(x^it^j)=
  \begin{cases}
  1/(j+1)& i=0\text{ and }j\text{ even},\\
  0&\text{otherwise}.
  \end{cases}
\]
This equals \(\CT_x(\tfrac12\int_{-1}^1 x^it^j\,dt)\). The least sphere
degree is the least total degree of a polynomial in \(x,y,t\)
representing the function on the sphere.

\begin{lemma}[Isolation of the top coefficient]\label{lem:isolation}
Suppose a regular phase \(f\) has highest weight \(n>0\), with part
\(t^n B(z)\) at that weight. Let \(q=p^r\) be an odd prime power and
\(m=(q-1)/n\) a positive integer. If the Laurent coefficients of \(f\)
are integral at a place over \(p\) of their number field, then
\(q\A(f^m)\) is integral there and
\[
                 q\A(f^m)\equiv \CT_z B(z)^m.
\]
The congruence is in the residue field of the place.
\end{lemma}

\begin{proof}
At angle zero the weight is the \(t\)-degree, which ranges from zero
to \(nm=q-1\) in \(f^m\). Of the allowed denominators only \(q\)
has valuation \(r\), and its degree \(q-1\) is even. The coefficient
there comes from the highest-weight part alone.
\end{proof}

We recall the cohomology calculation in the character convention used
below. Let \(Y^d=G(u)=\prod_b(u-b)^{a_b}\) define a geometrically
connected tame cyclic cover of \(\mathbb P^1\) after normalization, with
exponent \(a_\infty=-\sum_b a_b\bmod d\) at infinity. Write \(H_i\)
for the \(Y^i\)-character space in \(H^1(\mathcal O)\) over an algebraic
closure, for \(1\le i<d\). A generator of the character eigensheaf
on the finite chart is
\[
       w_i=Y^i\prod_b(u-b)^{-\lfloor ia_b/d\rfloor}.
\]
The finite factors divide out the largest integral powers of \(u-b\)
allowed by regularity. Normalizing at infinity gives the generator
\(u^{-b_i}w_i\), where \(b_i=\lceil\sum_b\{ia_b/d\}\rceil\).
It also generates on the finite part of the infinity chart
\(\mathbb P^1\setminus\{0\}\), since there \(u\) is a unit.
Thus the \v{C}ech quotient has basis \(w_i/u,\ldots,w_i/u^{b_i-1}\)
and its dimension is
\begin{equation}\label{eq:dimension}
       h(i)=-1+\sum_{b\,\text{including }\infty}\{ia_b/d\}.
\end{equation}
For nontrivial characters connectedness ensures a positive fractional
sum over all branches. If the exponents
are nonnegative and less than \(d\) on the finite chart and
\(d<\deg G<2d\), then \(H_1\) is one-dimensional with generator \(Y/u\).
When \(q=p^r\equiv1\pmod d\), the \(q\)-power map sends its representative
to \(YG(u)^{(q-1)/d}/u^q\). Hence the return coefficient on that line is
\begin{equation}\label{eq:return}
                   [u^{q-1}]G(u)^{(q-1)/d}.
\end{equation}

Absolute Frobenius acts on characters by \(i\mapsto pi\bmod d\),
since it raises both functions and eigenvalues to the \(p\)-th power.
The \(p\)-rank of the Jacobian is the dimension of the stable image
of this Frobenius on \(H^1(\mathcal O)\) \cite[\S2]{bouw}.
An orbit with unequal coherent dimensions is incompatible with an
invertible Frobenius. A zero space on an orbit forces zero stable rank
on that orbit. We write \(J^{[d]}\) for the factor of the Jacobian cut
out, up to isogeny, by the primitive order-\(d\) idempotent of the
deck group. It retains the spaces for \(\gcd(i,d)=1\).
At the good reductions considered below, this idempotent acts on
coherent cohomology since its denominator divides \(d\).
When a primitive
character sees exactly \(b\) nonzero branch exponents, its dimension
plus that of its inverse is \(b-2\), by \eqref{eq:dimension}.
For a quotient of degree prime to \(p\) that respects the order-\(d\)
action, the trace and pullback exhibit its coherent cohomology as a
Frobenius-stable summand. In characteristic zero the norm and pullback
likewise give an isogeny factor of the Jacobian.

\section{Fixed prime and unbounded conductor}

\begin{theorem}\label{thm:fixedprime}
Fix a rational prime \(p>3\) with \(p\equiv2\pmod3\) and an odd
integer \(r\ge3\). Set \(N=(p^r-1)/(p-1)\), and take \(c\in\Q\)
with \(v_3(c)=1\) and \(v_p(c)=v_p(1-c)=0\). The normalized
cover
\[
              C\colon W^N=z^{N-2}(z^2-1)(z^2-c)
\]
is connected. Its primitive Jacobian factor has dimension
\(2\varphi(N)\), is not potentially CM, and is nonordinary at \(p\).
At \(p\) the \(r\)-fold semilinear Frobenius fixes the basis class
\(W/z\) of its coherent character-one line.

For a regular sphere phase \(f\) of highest weight \(j+i=N\)
whose weight-\(N\) part is \(t^N z^{-2}(z^2-1)(z^2-c)\), with
\(z=x/t\), the least sphere degree is \(N\). The phase is not a
proper power over an algebraic closure of the coefficient field in
the sphere function field. For every such phase over a number field
with strict lower Laurent terms integral at a place over \(p\),
for every integer \(k\ge1\) we have
\[
         p^{kr}\A(f^m)\equiv1,\qquad m=\frac{p^{kr}-1}{N},
\]
in the residue field. In particular the \((p-1)\)-st moment is nonzero.
\end{theorem}

\begin{proof}
We have \(\gcd(N,6p)=1\) because \(p\equiv2\pmod3\) and \(r\) is
odd. The exponent one at \(z=1\) gives geometric connectedness.
With \(u=z^2\) the function field also has the model
\[
          W^{2N}=u^{N-2}(u-1)^2(u-c)^2.
\]
Indeed \(z^2\) and \(z^{N-2}\) recover \(z\) because \(N\) is odd.
The original deck group is generated by the square of a generator of
the refined group of order \(2N\). With \(V=W^2\) the even refined
characters are the pullback from its degree-two quotient
\[
          D\colon V^N=u^{N-2}(u-1)^2(u-c)^2.
\]
There are six branch points including infinity on \(C\) over \(z\),
and four on \(D\) over \(u\), all seen by each primitive character.
Their primitive dimensions are \(2\varphi(N)\) and \(\varphi(N)\).

\smallskip
\noindent\textit{Toric reduction at \(3\).}
Blow up the collision of \(0,c\) in the base of \(D\). Since
\(v_3(c)=1\), the strict transforms of the branch sections separate,
and at the node the equation of the base is \(us=c\). The Kummer
equation there is
\[
                (V/u)^N=(u-1)^2(1-s)^2.
\]
The right-hand side is a unit at the node. Since \(3\nmid N\),
the normalization is \'{e}tale over the node with \(N\) geometric
nodes, permuted regularly by the deck group. The main component has
equation \((V/u)^N=(u-1)^2\). On the exceptional component, using
\(U=u/c\), the equation is \((V/c)^N=U^{N-2}(U-1)^2\). Both
components are connected rational curves, as each is a cyclic cover
with two fully ramified points. Away from the nodes, the normalization
is smooth at the separated tame branch sections. In local parameters
a unit times a single parameter power gives the usual tame cyclic
normalization, smooth after adjoining the needed unit roots over a
strict henselization.

The torus of the Jacobian of a semistable curve has character lattice
the first homology of its dual graph \cite[(12.3.7), \S12.4]{sga}.
Here the graph has two vertices joined by \(N\) edges, giving the
augmentation representation of the deck group. Its primitive part
has rank \(\varphi(N)\), equal to the dimension of \(J^{[N]}(D)\).
This toric rank persists after finite extension. CM abelian varieties
have potentially good reduction everywhere, and potential good
reduction is inherited by isogeny factors by the inertia criterion
\cite[Theorems~1 and~6]{serretate}. The positive-dimensional toric
factor of \(D\) therefore precludes potential CM for \(J^{[N]}(C)\).

\smallskip
\noindent\textit{A failing orbit and an invertible return at \(p\).}
The branch points remain distinct and tame at \(p\). For a primitive
original label \(i\), write \(x=i/N\) using the representative in
\([1,N-1]\). Its two refined labels have parities \(e=0,1\).
The branch exponents \((N-2,2,2,N-2)\) of the refined model give
coherent dimension
\[
       -1+2\{e/2-x\}+2x
          =\begin{cases}1,&e=0,\\2\mathbf 1_{x>1/2},&e=1.\end{cases}
\]
Indeed the two labels satisfy \(\ell\equiv i\pmod N\), so the
fractional parts of \(\ell(N-2)/(2N)\) and \(2\ell/(2N)\)
are \(\{e/2-x\}\) and \(x\). Frobenius preserves the parity.
The odd refined label over \(i=2\) has dimension zero. With
\(X=p^{r-1}\), its iterate over \(p^{r-1}i\bmod N\) has
representative \(2X-N\), with
\[
                 \frac{2X-N}{N}=\frac{(p-2)X+1}{pX-1}
                      \in(1/2,1).
\]
Its coherent dimension is two. This orbit has zero stable rank and
proves nonordinarity.

On the original character-one line the generator is \(W/z\), as
the finite exponents are less than \(N\) and sum to \(N+2\).
Put \(m=p-1\). The scalar \eqref{eq:return} at \(q=p^r\) is
\begin{align*}
 [z^{q-1}]\bigl(z^{N-2}(z^2-1)(z^2-c)\bigr)^m
     &=[u^m](u-1)^m(u-c)^m\\
     &=\sum_{a=0}^{p-1}c^a=(c-1)^{p-1}=1\quad\text{in }\F_p.
\end{align*}
The coefficient identity uses
\((u-b)^{p-1}=\sum_{a=0}^{p-1}b^{p-1-a}u^a\) in characteristic
\(p\). The residue of \(c\) lies in \(\F_p\) and differs from one.
The \(r\)-fold Frobenius fixes the basis class, so its \(k\)-fold
iterate does as well. For each \(q=p^{kr}\), put \(m=(q-1)/N\).
The return coefficient \eqref{eq:return} is again the angular
constant term \(\CT_z(z^{-2}(z^2-1)(z^2-c))^m\).
Lemma~\ref{lem:isolation} gives all the moment congruences.

\smallskip
\noindent\textit{Degree and powers.}
The sphere relation eliminates every occurrence of \(xy\) without
raising total degree. The remaining monomials in \(x,t\) and in
\(y,t\) with positive \(y\)-exponent form a basis: distinct angles
are independent, and within angle \(-v<0\) the Laurent monomials
obtained from \(y^vt^b\) have distinct highest \(t\)-degrees \(2v+b\).
Thus the highest weight \(j+i\) is exactly the highest degree in
this normal form. One compatible lift is
\[
         4c y^2t^{N-2}-(1+c)t^N+x^2t^{N-2}.
\]
Finally, in the function field with coordinates \(z,t\), the pole
order of \(f\) at \(t=\infty\) is \(N\). A proper power of exponent
\(d>1\) would have \(d\mid N\) and its leading coefficient at
\(t=\infty\) would be a \(d\)-th power in the residue field. The
coefficient has pole order two at \(z=0\), impossible for such \(d\)
since \(N\) is odd.
\end{proof}

\begin{proposition}[Exact \(5\)-rank]\label{prop:rankhalf}
In Theorem~\ref{thm:fixedprime} take \(p=5\) and \(c=3\). Let \(S_r\)
consist of the integers \(1\le i<N\) coprime to \(N\) for which all
\(r\) base-\(5\) digits of \(4i\), including leading zeros, are at
least two. Then the \(5\)-rank of \(J^{[N]}(C)\) is
\[
                      \varphi(N)+2\lvert S_r\rvert.
\]
Its ratio to the dimension tends to \(1/2\) as odd \(r\) tends to
infinity. At \(N=31\) this gives \(5\)-rank \(42\).
\end{proposition}

\begin{proof}
Work on the refined cover of order \(2N\), where parity is preserved
by Frobenius. Let \(w_\ell\) be the finite generator for a refined
character of parity \(e\), and let \(\ell'\equiv5\ell\pmod{2N}\).
Let \(i\) be the least positive residue of \(\ell\) modulo \(N\).
Write \(x=i/N\), \(a=\lfloor5x\rfloor\), and \(x'=5x-a\). In
reduction modulo \(5\), put \(Q(u)=(u-1)(u-3)=u^2+u+3\).
The finite generator formula gives
\(w_\ell^5=w_{\ell'}u^gQ(u)^a\), where
\(g=\lfloor5\{e/2-x\}\rfloor\) is the exponent at zero.
Indeed, each exponent in the extra factor is the floor of five
times the corresponding fractional part at the source.
The entry from \(w_\ell/u^j\) to \(w_{\ell'}/u^{j'}\) is therefore
\([u^{5j-j'-g}]Q(u)^a\).

For even parity \(g=4-a\). The line-to-line entries for
\(a=0,1,2,3,4\) are \([u^a]Q(u)^a=1,1,2,4,1\) in \(\F_5\),
so the even part contributes \(\varphi(N)\).
For odd parity a Frobenius orbit with any \(x<1/2\) contributes zero
by the dimension formula in the preceding proof. A step between two
spaces of dimension two has \(x,x'>1/2\) and \(g=6-a\).
For \(a=2,3,4\) the matrices on the indicated \v{C}ech bases are,
respectively,
\[
 \begin{pmatrix}4&0\\0&1\end{pmatrix},\qquad
 \begin{pmatrix}2&1\\2&3\end{pmatrix},\qquad
 \begin{pmatrix}2&4\\3&3\end{pmatrix}.
\]
All have determinant \(4\) in \(\F_5\). Hence an odd orbit with
every label above \(1/2\) contributes two per label.

Now \(x=4i/(5^r-1)\) has periodic base-\(5\) digits given by the
\(r\)-digit block of \(4i\). Frobenius cyclically rotates that block.
The expansion of \(1/2\) consists entirely of twos. A digit less
than two forces some rotation below \(1/2\), and a block of digits
at least two has every rotation strictly above \(1/2\) because
\(N\) is odd and \(x\ne1/2\). This proves the formula. We have
\(\lvert S_r\rvert\le3^r\) and \(\varphi(N)\ge N^{3/4}\), since
each prime \(\ell\mid N\) is at least five and
\(1-1/\ell\ge\ell^{-1/4}\). As \(5^{3/4}>3\), the extra term is
\(o(\varphi(N))\). For \(r=3\), the allowed blocks are the six
permutations of \((2,2,4)\) and \((2,3,3)\), by divisibility of
\(4i\) by four and \(0<4i<5^3-1\).
\end{proof}

\section{A fixed cover and its certifying prime classes}

\begin{theorem}\label{thm:fixedcover}
Let \(M\ge5\) with \(M\equiv5\pmod{14}\), put \(n=7M\), and set
\[
           P(U)=(U-1)^{M+1}(U-c)^{6M},\qquad c\in\Q,\quad v_2(c)=1.
\]
The normalized cover \(C\colon Y^n=P(z^7)\) is connected and its
primitive factor has dimension \(13\varphi(n)/2\). This factor is
not potentially CM. At every prime \(p\nmid n\) with
\(v_p(c)=v_p(1-c)=0\), it is nonordinary if \(p\not\equiv1\pmod n\).

Let \(f\) be a fixed regular phase over a characteristic-zero field
of highest weight \(j+2i=n\), with weight-\(n\) part
\(t^n z^{-7n}P(z^7)\), where \(z=x/t^2\). For a prime \(p\) as above,
let \(d\) be its order modulo \(n\). Write \(\rho(i)\) for the
least positive residue of \(i\) modulo seven. The certifying criterion is
\begin{equation}\label{eq:orbitcriterion}
           \frac{i}{n}+\frac{\rho(i)}7<1
         \qquad\text{for each least positive representative }i\in\langle p\rangle.
\end{equation}
If the phase is defined over a number field and integral at a place
over \(p\), then for every \(k\ge1\) the normalized moment
\(p^{kd}\A(f^{(p^{kd}-1)/n})\) has nonzero residue if and only if
\eqref{eq:orbitcriterion} holds.
For any fixed \(f\), all sufficiently large primes meeting the criterion
certify nonzero moments at these orders. The nonsplit certifying classes
have natural density at least \(5/\varphi(n)\), and exactly \(7/24\) when
\(n=35\).

All compatible strict lower terms are allowed. Such phases are not proper
powers in the sphere function field over an algebraic closure, and there
is a compatible phase of least sphere degree \(8n\).
\end{theorem}

\begin{proof}
The exponent \(M+1\) at a root of \(z^7=1\) is coprime to \(n\),
giving connectedness. There are fourteen finite branch points and
one at infinity, whose exponent is \(-7\bmod n\). Each primitive
character sees all fifteen, giving dimension \(13\varphi(n)/2\)
by \eqref{eq:dimension}. Over \(U=z^7\) there are two commuting deck
actions, of orders \(n\) and seven. For \(0\le k\le6\), the functions
\(U\) and \(Z_k=Yz^k\) give a quotient of degree seven, with model
\[
             D_k\colon Z_k^n=U^{kM}P(U).
\]
Indeed the two functions are fixed by the order-seven action
\((z,Y)\mapsto(\xi z,\xi^{-k}Y)\) for \(\xi^7=1\), and adjoining
\(z\) recovers the original field. The action on \(Y\) descends to
the action on \(Z_k\), so the primitive part of each quotient
pulls back to an isogeny factor of \(J^{[n]}(C)\).

\smallskip
\noindent\textit{The toric quotient.}
In \(D_1\) the colliding branch exponents at \(2\), at \(U=0,c\),
are \(M\) and \(6M\). After blowing up the base, the local equation
\(Us=c\) at the node gives
\[
        (Z_1/U)^n=(U-1)^{M+1}(1-s)^{6M}.
\]
The normalization is \'{e}tale at the node with \(n\) geometric
nodes. The main component is the \(n\)-th root of \((U-1)^{M+1}\)
and is connected and rational. On the exceptional component, use
\(u=U/c\) and \(Z_1/c\). The reduced right-hand side is a nonzero
constant times \(u^M(u-1)^{6M}\). This has \(M\) geometric components,
each a degree-seven cyclic cover with two branch points, hence
rational. Away from the node the same tame normalization argument as
for \(D\) in Theorem~\ref{thm:fixedprime} gives smoothness.

The vertex representation of the dual graph has one fixed vertex
and \(M\) permuted vertices, on which the deck group acts through a
quotient of order \(M\). It contains no primitive order-\(n\)
characters. The \(n\) edges carry the regular representation, so the
primitive part of graph homology has dimension \(\varphi(n)\).
The primitive factor of \(D_1\) has precisely this dimension, since
it has four branch points seen by each primitive character. As
above, this totally toric factor proves that \(J^{[n]}(C)\) is not
potentially CM.

\smallskip
\noindent\textit{Nonordinarity at nonsplit primes.}
The branch points remain distinct and tame at the primes under
consideration. For a primitive label \(i\), write \(x=i/n\),
\(t=\rho(i)\), and \(v=(kt\bmod7)/7\) with \(0\le v<1\).
The four exponents of \(D_k\) give the fractions of \(v,x+t/7,
1-t/7,-x-v\). Writing \(\epsilon_j(x)=\mathbf 1_{x+j/7>1}\),
the dimension formula \eqref{eq:dimension} gives
\[
 h_k(i)=1-\epsilon_t(x)+\epsilon_{kt\bmod7}(x),\qquad
 \sum_{k=0}^6\bigl(h_k(i)-h_0(i)\bigr)=\left\lfloor\frac{i}{M}\right\rfloor.
\]
There is no equality \(x+j/7=1\) because \(M>1\) and \(i\) is
primitive. The sum counts the residues \(j=0,\ldots,6\) with
\(x+j/7>1\).

Ordinarity of the primitive factor of \(C\) would make all the quotient
factors ordinary, forcing \(h_k(i)\) constant on every Frobenius
orbit. Consequently multiplication by \(a=p\bmod n\) would preserve
\(\lfloor i/M\rfloor\) for all primitive labels. Testing \(i=1\)
gives \(1\le a<M\). If \(a>1\), choose the first power of two \(j\)
with \(aj\ge M\). Then \(j<M\) and \(M\le aj<2M\). Both are
primitive since \(n\) is odd and \(a\) is a unit. The two floors
differ, a contradiction. Hence \(a=1\) is necessary for ordinarity.

\smallskip
\noindent\textit{The certifying criterion.}
For \(D_0\), the polynomial \(P\) has degree \(n+1\), and its
two finite exponents are less than \(n\). Its character-one line
has generator \(Y/U\), which pulls back to \(Y/z^7\) on \(C\).
For \(q=p^d\) and \(m=(q-1)/n\), the return coefficient
\eqref{eq:return} is \(\lambda_p=[U^{q-1}]P(U)^m\in\F_p\).
It is also the angular constant term in Lemma~\ref{lem:isolation},
since \(\CT_z z^{-7nm}P(z^7)^m=[U^{q-1}]P(U)^m\).

In characteristic \(p\), translation by one preserves the coefficient
at degree \(q-1\) of a polynomial of degree at most \(2q-2\).
For every term of degree \(q\le b\le2q-2\), the factorization
\((V+1)^b=(V^q+1)(V+1)^{b-q}\) gives zero coefficient at \(q-1\).
Since \(\deg P^m=q-1+m\le2q-2\), translate \(U=V+1\) to obtain
\begin{equation}\label{eq:binomial}
         [U^{q-1}]P(U)^m
             =(1-c)^m\binom{6Mm}{m}\qquad\text{in }\F_p.
\end{equation}
In fact \(P(V+1)^m=V^{(M+1)m}(V+1-c)^{6Mm}\), and the required
degree in the last factor is \(6Mm-m\).

The \(d\) digits of \(m\) and \(6Mm\) come from long division of
\(1/n\) and \(6/7\). For each successive \(i=p^j\bmod n\) the
two normalized remainders are \(i/n\) and \(1-\rho(i)/7\).
Under \eqref{eq:orbitcriterion}, taking floors after multiplication
by \(p\) shows that each digit of \(6Mm\) is at least the
corresponding digit of \(m\). The base-\(p\) factorization of
\((1+X)^{6Mm}\) makes \eqref{eq:binomial} a unit. Conversely, if
\eqref{eq:orbitcriterion} fails somewhere, then \(h_0(i)=0\) there,
so the return on the character-one line passes through zero.
The \(kd\)-fold return is multiplication of the basis class by
\(\lambda_p^k\), since \(\lambda_p\in\F_p\), and
Lemma~\ref{lem:isolation} gives that same residue at order
\((p^{kd}-1)/n\). This proves the criterion for the residue.

For a fixed phase over a number field, the coefficients are integral
away from finitely many primes.
For any characteristic-zero field, take the finitely generated
\(\Q\)-algebra generated by the coefficients of the polynomial normal
form of \(f\). A closed point gives a regular phase over a number field
with the same rational leading form, since the passage to Laurent
coefficients is rational linear.
Each moment is a rational polynomial in the finitely many Laurent
coefficients of the phase and so commutes with this specialization.
A nonzero specialized moment certifies the original at each order
after discarding the finitely many bad primes of the specialized phase.

\smallskip
\noindent\textit{Density.}
The class \(M-1\) generates a subgroup of order six, since it is
\(-1\bmod M\) and \(4\bmod7\). The least representatives of its
powers are
\[
          1,\ M-1,\ 1+3M,\ 6M-1,\ 1+2M,\ 2M-1.
\]
All satisfy \(i+M\rho(i)<n\), so the five nonidentity classes
certify at nonordinary primes. Testing \eqref{eq:orbitcriterion}
on the \(24\) units modulo \(35\) gives \(1,4,8,9,11,16,22,29\).
The prime number theorem in arithmetic progressions gives the
claimed densities.

\smallskip
\noindent\textit{Sphere lifts.}
The angular exponents in the leading form range from \(-7n\) to
seven. At a negative angle \(-v\), lift the prescribed coefficient
times \(t^n z^{-v}\) by multiplying that coefficient by
\((-2)^v y^v t^n\). At angles zero and seven use \(t^n\) and
\(x^7t^{n-14}\), respectively, with the prescribed coefficients.
These choices have the indicated leading form because \(2xy=1-t^2\).
The lift contains \(y^{7n}t^n\) with nonzero coefficient, so its
normal-form degree and least sphere degree are \(8n\).
For any phase with this leading form, a proper power would have
exponent dividing the pole order \(n\) at \(t=\infty\) in the
coordinates \(z,t\). The leading angular coefficient has order
\(M+1\) at a root of \(z^7=1\), coprime to \(n\), excluding such a
power.
\end{proof}

\begin{samepage}
\paragraph{AI assistance and review.}
ChatGPT (OpenAI) assisted with mathematical arguments, literature checks,
exposition, proof review and verification code.

\begin{thebibliography}{9}
\bibitem{bouw}
I. I. Bouw, The \(p\)-rank of ramified covers of curves,
\emph{Compos. Math.} \textbf{126} (2001), 295--322.
\href{https://doi.org/10.1023/A:1017513122376}{DOI: 10.1023/A:1017513122376}.

\bibitem{lmpt}
W. Li, E. Mantovan, R. Pries and Y. Tang, Newton polygons of cyclic
covers of the projective line branched at three points, in \emph{Research
Directions in Number Theory}, Assoc. Women Math. Ser. \textbf{19},
Springer, 2019, 115--132.
\href{https://arxiv.org/abs/1805.04598}{arXiv:1805.04598}.

\bibitem{sga}
A. Grothendieck, Mod\`eles de N\'eron et monodromie,
in \emph{Groupes de monodromie en g\'eom\'etrie alg\'ebrique I},
S\'eminaire de G\'eom\'etrie Alg\'ebrique du Bois-Marie 1967--1969 (SGA 7 I),
Lecture Notes in Math. \textbf{288}, Springer, 1972, Expos\'e IX.
\href{https://grothendiecksga.com/read/sga7/en/9-12.html}{\S12}.

\bibitem{serretate}
J.-P. Serre and J. Tate, Good reduction of abelian varieties,
\emph{Ann. of Math.} \textbf{88} (1968), 492--517.
\href{https://doi.org/10.2307/1970722}{DOI: 10.2307/1970722}.
\end{thebibliography}
\end{samepage}
\end{document}
